% TOP_PROBLEM: {"id": 3000004, "title": "Are t-perfect graphs strongly t-perfect?", "category_id": 2, "category": {"id": 2, "name": "combinatorics", "display_name": "Combinatorics", "description": "Counting problems, graph theory, discrete structures.", "slug": "combinatorics", "order_index": 2, "created_at": "2026-07-31T15:26:25.671Z"}, "status": "open", "problem_number": "AMR-029-0004", "set_id": 13}
% FIRST_POSTED: 2026-10-01
% ENTRY_KIND: statement_only
% UnsolvedMath ID 3000004; source catalogue code AMR-029-0004
% !TeX program = lualatex
% Definition and sources only; this file is not an LLM solution attempt.
\documentclass[11pt,a4paper]{article}
\usepackage[margin=25mm]{geometry}
\usepackage{fontspec}
\setmainfont{lmroman10-regular.otf}[BoldFont=lmroman10-bold.otf,
 ItalicFont=lmroman10-italic.otf,BoldItalicFont=lmroman10-bolditalic.otf]
\usepackage{amsmath,amssymb,mathtools}
\usepackage{unicode-math}
\setmathfont{latinmodern-math.otf}
\AtBeginDocument{\let\setminus\smallsetminus}
\usepackage{xurl,hyperref}
\hypersetup{colorlinks=true,urlcolor=blue}
\setcounter{secnumdepth}{0}
\setlength{\parindent}{0pt}
\setlength{\parskip}{5pt}
\setlength{\emergencystretch}{3em}
\begin{document}
\section*{Are t-perfect graphs strongly t-perfect?}\label{3000004}
{\small UnsolvedMath ID 3000004 · AMR-029-0004 · Combinatorics · AMR Open Problem Lists}
\subsection{Definitions and mathematical statement}
For a finite simple graph $G=(V,E)$, a stable set contains no adjacent
pair. Its stable-set polytope is the convex hull of the vectors
$\mathbf1_I\in\{0,1\}^V$ of stable sets $I$. For $U\subseteq V$, write
$x(U)=\sum_{v\in U}x_v$. Consider the system
\[
 0\le x_v\le1\ (v\in V),\qquad x_u+x_v\le1\ (uv\in E),
 \qquad x(V(C))\le (|V(C)|-1)/2
\]
for every odd simple cycle $C$ of $G$. Call its feasible polytope $Q(G)$.
The graph is $t$-perfect when $Q(G)$ equals its stable-set polytope.

For a rational system $Ax\le b$, total dual integrality means that, for
every integral objective vector $c$ whose primal maximum is finite, the
dual program
\[
             \min\{b^Ty:A^Ty=c,\ y\ge0\}
\]
has an integral optimal solution. A graph is strongly $t$-perfect when
the displayed defining system for $Q(G)$ is totally dual integral.
Does $t$-perfection always imply strong $t$-perfection?

The defining rows and their normalization matter: total dual integrality
is a property of a system, not merely of its feasible set. Vertex upper
bounds explicitly handle isolated vertices; for other vertices they follow
from an incident edge inequality and nonnegativity. The question asks for
integral dual certificates for every integral objective, not just integral
vertices of the primal polytope.
\subsection{Short English statement}
Does the usual odd-cycle description of a $t$-perfect graph's
stable-set polytope always provide integral optimal dual certificates?
\subsection{Sources}
\begin{itemize}
\item[S1] ULAM.AI and the UnsolvedMath Contributors, supplied problems.json, record 3000004 (AMR-029-0004), AMR Open Problem Lists. Catalogue identity and source transcription; CC BY 4.0.
\url{https://huggingface.co/datasets/ulamai/UnsolvedMath}.
\item[S2] Egres Open - List of open problems, Open-problem page: Are t-perfect graphs strongly t-perfect?. Original source indexed by AMR-029-0004.
\url{https://oldlemon.cs.elte.hu/egres/open/List_of_open_problems}.
\item[S3] EGRES Open, Are t-perfect graphs strongly t-perfect?. Individual source page and explanatory remarks.
\url{https://lemon.cs.elte.hu/egres/open/Are_t-perfect_graphs_strongly_t-perfect%3F}.
\end{itemize}
\end{document}
